\section{Radial channels: bypass, local synthesis, and gain limits}
\label{cn:section}

The large multiplicity $2\ell+1$ of a spherical partial wave and the
unbounded radial Bragg amplification in Proposition~\ref{cx:bragg}
suggest a concrete possible counterexample.  This section tests that
proposal directly at the physical section.  The resulting obstruction
does not assume a field-energy bound or a bound on the exterior support.
It applies to a single angular order in a homogeneous core, including
arbitrary coherent combinations of all azimuthal indices and both
electric and magnetic multipoles.  Mixed angular orders and a structured
core remain outside its scope.

\subsection{Inversion parity and the threshold regime}

\begin{proposition}[PROVED: an odd plane flux must return outside the section]
\label{cn:odd-return}
Let $s\in L^1(\mathbb R^2)$ be a real normal power density, and suppose
$s(2c-x)=-s(x)$ almost everywhere for some point $c$ in the section
plane.  For any measurable section $D$,
\begin{equation}\label{cn:odd-return-bound}
 \int_{\mathbb R^2\setminus D}|s|\,dA
 \ \geq\ \left|\int_Ds\,dA\right|.
\end{equation}
Consequently a nonzero counted input with this symmetry requires
$\delta\geq\eta\tau$.  This conclusion does not require that $D$ be
centered at $c$.
\end{proposition}
\begin{proof}
The intersection $D\cap(2c-D)$ is invariant under $x\mapsto2c-x$,
so its signed integral is zero.  The reflection of
$D\setminus(2c-D)$ is contained in the complement of $D$, and has the
opposite signed integral.  Taking its absolute integral proves
\eqref{cn:odd-return-bound}.  The counted-input inequalities then give
$\delta\|a\|^2\geq P_D(a)\geq\eta\tau\|a\|^2$.
\end{proof}

\begin{remark}[PROVED: application to a Maxwell parity sector]
\label{cn:parity}
In a medium invariant under inversion about $c$, a field satisfying
$E(2c-x)=\sigma E(x)$, $\sigma\in\{1,-1\}$, satisfies
$H(2c-x)=-\sigma H(x)$ by the first curl equation.  Its real Poynting
vector is therefore odd.  If $c$ lies in the section plane and the
required traces are integrable, Proposition~\ref{cn:odd-return} applies.
In particular, a coherent sum of electric-type multipoles of one
angular order has one parity, as does a coherent sum of magnetic-type
multipoles of that order.  Merely using their azimuthal degeneracy does
not give passage in the regime $\delta<\eta\tau$.  If the square is
centered at $c$, its signed passage is exactly zero.  Mixing the two
parities is necessary for this centered construction; it is not
sufficient, as the next theorem shows.
\end{remark}

\subsection{Both polarizations of a regular spherical partial wave}

Fix a homogeneous ball centered at the center of $D$, with real
permittivity $n^2\eps_0$ and permeability $\mu_0$.  Write
$k=nk_0$ and $Z=\sqrt{\mu_0/(n^2\eps_0)}$.  Let
$\nu=\ell(\ell+1)$, $\ell\geq1$, and let $Y_{\ell m}$ be any
orthonormal complex spherical-harmonic basis.  For coefficient vectors
$A,B\in\mathbb C^{2\ell+1}$ define
\begin{equation}\label{cn:angular-fields}
 Y_A=\sum_m A_mY_{\ell m},\quad
 \Psi_A=\nu^{-1/2}\nabla_{\mathbb S^2}Y_A,\quad
 X_A=\widehat r\times\Psi_A,
\end{equation}
and similarly for $B$.  With $j=j_\ell(kr)$, put
\begin{equation}\label{cn:radial-fields}
 u=rj,\qquad a_\ell(r)=\frac{\sqrt\nu\,j}{kr},\qquad
 b_\ell(r)=\frac{u'}{kr}.
\end{equation}
The notation $a_\ell,b_\ell$ in this section denotes radial functions,
not external incoming amplitudes.  The general regular field at this
one angular order can be written
\begin{align}
 E&=jX_A-a_\ell Y_B\widehat r-b_\ell\Psi_B,\notag\\
 H&=-\frac{i}{Z}
       (-a_\ell Y_A\widehat r-b_\ell\Psi_A+jX_B).
 \label{cn:multipoles}
\end{align}
The usual two regular spherical Maxwell waves give this representation:
if $M_A=jX_A$ and $N_A=k^{-1}\nabla\times M_A$, direct spherical
differentiation gives
$N_A=-a_\ell Y_A\widehat r-b_\ell\Psi_A$,
$\nabla\times N_A=kM_A$, and hence both Maxwell curl equations
hold in \eqref{cn:multipoles}.  Regularity excludes the singular
spherical Bessel solution.  All possible complex coefficients,
including those generated by a distant resonant dielectric, are allowed.

\begin{lemma}[PROVED: the plane flux separates into opposite parities]
\label{cn:flux-split}
On the equatorial plane $z=0$, a field \eqref{cn:multipoles} has
\begin{equation}\label{cn:flux-decomposition}
 S_z(r,\varphi)=F_\ell(r)O(\varphi)+G_\ell(r)C(\varphi),
\end{equation}
where the real angular functions satisfy
$O(\varphi+\pi)=-O(\varphi)$ and
$C(\varphi+\pi)=C(\varphi)$, and
\begin{equation}\label{cn:radial-flux-factors}
 F_\ell=\frac{a_\ell j}{2Z}
        =\frac{\sqrt\nu\,u^2}{2Zkr^3},\qquad
 G_\ell=\frac{a_\ell b_\ell}{2Z}
        =\frac{\sqrt\nu\,u u'}{2Zk^2r^3}.
\end{equation}
Explicitly, with all angular expressions evaluated at $\theta=\pi/2$,
\begin{align}
 O&=-\operatorname{Im}
       (\overline{Y_A}\Psi_{A,\theta}
         +\overline{Y_B}\Psi_{B,\theta}),\notag\\
 C&=\operatorname{Im}
       (Y_B\overline{X_{A,\theta}}
         -\overline{Y_A}X_{B,\theta}).
 \label{cn:angular-flux-factors}
\end{align}
\end{lemma}
\begin{proof}
At the equator $\widehat z=-\widehat\theta$.  In
$E\times\overline H$, the tangential terms arise only by crossing one
radial and one tangential component.  Substituting
\eqref{cn:multipoles}, and using
$X\times\widehat r=\Psi$,
$\Psi\times\widehat r=-X$,
$\widehat r\times X=-\Psi$, gives
\eqref{cn:flux-decomposition}--\eqref{cn:angular-flux-factors}.
For completeness, the tangential vector inside the factor $i/Z$ is
\begin{equation}\label{cn:cross-product-check}
 a_\ell j(Y_B\overline{\Psi_B}-\overline{Y_A}\Psi_A)
 +a_\ell b_\ell
       (Y_B\overline{X_A}-\overline{Y_A}X_B).
\end{equation}
Taking one half its real part after multiplication by $i/Z$, then
changing from the $\theta$ component to $z$, gives the stated signs.

The spherical-harmonic inversion identity gives
$Y_A(\pi/2,\varphi+\pi)=(-1)^\ell Y_A(\pi/2,\varphi)$,
whereas its $\theta$ derivative gains the opposite sign.  Since
$X_\theta=-\nu^{-1/2}\partial_\varphi Y$ at the equator,
$X_\theta$ gains the same sign as $Y$.  These facts prove the two
angular parity assertions for arbitrary coherent coefficients.
\end{proof}

\begin{lemma}[PROVED: radial growth below the angular turning scale]
\label{cn:radial-growth}
Let $R_2>0$ and assume $\nu>k^2R_2^2$.  Define
\begin{equation}\label{cn:exponent}
 \gamma=\frac{1+\sqrt{1+4(\nu-k^2R_2^2)}}2>1.
\end{equation}
Then the regular real solution $u=rj_\ell(kr)$ is positive with
$u'>0$ for $0<r\leq R_2$, and
\begin{equation}\label{cn:log-derivative}
 \gamma\leq\frac{ru'(r)}{u(r)}\leq\ell+1,
 \qquad
 \frac{u(t)}{u(s)}\geq(t/s)^\gamma
 \quad(0<s\leq t\leq R_2).
\end{equation}
\end{lemma}
\begin{proof}
The radial equation is
$u''=(\nu/r^2-k^2)u$ and
$u(r)=c r^{\ell+1}(1+O(r^2))$ with $c>0$.
The coefficient of $u$ is positive throughout this interval, so a
first zero of $u$ or $u'$ is impossible.  For $w=ru'/u$ the equation is
\begin{equation}\label{cn:riccati}
 rw'=\nu-k^2r^2+w-w^2.
\end{equation}
At $w=\gamma$ its right side is
$k^2(R_2^2-r^2)\geq0$, and at $w=\ell+1$ it is $-k^2r^2<0$.
The initial expansion has
$w=\ell+1-k^2r^2/(2\ell+3)+O(r^4)$, so $w$ begins in the
interval $[\gamma,\ell+1]$ and cannot exit it.  A first crossing at
the lower endpoint can equivalently be excluded by the scalar
comparison equation, including the equality case at $R_2$.
Integrating $u'/u\geq\gamma/r$ proves the ratio bound.
\end{proof}

\begin{theorem}[PROVED: high angular degeneracy cannot evade absolute bypass in a homogeneous core]
\label{cn:radial-bypass}
Let $R=L/\sqrt2$ be the circumradius of the fixed square $D$, fix
$1<s<t$, and assume a homogeneous source-free ball of radius $tR$
about its center.  Consider any regular Maxwell field in this ball
supported in one angular order $\ell$, with arbitrary coefficients
$A,B$ in \eqref{cn:multipoles}.  If
$\ell(\ell+1)>k^2t^2R^2$, let $\gamma$ be
\eqref{cn:exponent} with $R_2=tR$, and set
\begin{equation}\label{cn:bypass-constant}
 B_\ell(s,t)=\frac{\gamma}{\ell+1}
       \bigl(t^{2\gamma-2}-s^{2\gamma-2}\bigr).
\end{equation}
Then its physical flux obeys the local, amplitude-independent estimate
\begin{equation}\label{cn:bypass-bound}
 \int_{sR<r<tR,\ z=0}|S_z|\,dA
 \ \geq\ B_\ell(s,t)\,|P_D|.
\end{equation}
In particular any nonzero counted input whose core field is of this
angular order must satisfy
\begin{equation}\label{cn:counted-obstruction}
 \delta\geq B_\ell(s,t)\eta\tau.
\end{equation}
At fixed $k,R,s,t$, $B_\ell(s,t)\to\infty$ exponentially as
$\ell\to\infty$.  Neither exterior field energy nor a radial mirror's
quality factor appears in this inequality.
\end{theorem}
\begin{proof}
By Lemma~\ref{cn:radial-growth}, $G_\ell>0$ on $(0,tR]$.
The square is centrally symmetric, so the integral of the odd term
$F_\ell O$ in \eqref{cn:flux-decomposition} over $D$ is zero.
Writing the square in polar coordinates and using its circumradius,
\begin{equation}\label{cn:inside-angular-bound}
 |P_D|\leq
 \left(\int_0^R G_\ell(r)r\,dr\right)
 \left(\int_0^{2\pi}|C(\varphi)|\,d\varphi\right).
\end{equation}
On the full annulus, opposite angles have fluxes $FO+GC$ and
$-FO+GC$.  The triangle inequality
$|x+y|+|-x+y|\geq2|y|$ gives
\begin{equation}\label{cn:outside-angular-bound}
 \int_{sR<r<tR}|S_z|\,dA\geq
 \left(\int_{sR}^{tR}G_\ell(r)r\,dr\right)
 \left(\int_0^{2\pi}|C(\varphi)|\,d\varphi\right).
\end{equation}
There is no basis-only estimate here: both bounds apply to the full
coherent field with its actual $A,B$.

The common positive factor $\sqrt\nu/(2Zk^2)$ cancels in the ratio
of the radial integrals, leaving integrals of $u u'/r^2$.
By \eqref{cn:log-derivative},
\begin{align}
 \int_0^R\frac{u u'}{r^2}\,dr
 &\leq\frac{\ell+1}{2\gamma-2}\frac{u(R)^2}{R^2},\notag\\
 \int_{sR}^{tR}\frac{u u'}{r^2}\,dr
 &\geq\frac{\gamma}{2\gamma-2}\frac{u(R)^2}{R^2}
       \bigl(t^{2\gamma-2}-s^{2\gamma-2}\bigr).
 \label{cn:radial-integral-ratio}
\end{align}
For the first inequality use
$u(r)\leq u(R)(r/R)^\gamma$ and $u'\leq(\ell+1)u/r$;
for the second use the reverse growth bound for $r\geq R$ and
$u'\geq\gamma u/r$.  Combining
\eqref{cn:inside-angular-bound}--\eqref{cn:radial-integral-ratio}
proves \eqref{cn:bypass-bound}; if $C$ vanishes identically, then
$P_D=0$ and the assertion is immediate.  The annulus lies outside
$D$, so the counted-input conditions imply
\eqref{cn:counted-obstruction}.  Finally
$\gamma=\ell+1+O(\ell^{-1})$ at fixed $ktR$, which proves the
claimed divergence.
\end{proof}

\begin{corollary}[PROVED: a limitation of the degenerate Bragg proposal]
\label{cn:bragg-consequence}
Fix the homogeneous core, the section, and the annulus in
Theorem~\ref{cn:radial-bypass}, together with finite
$\delta\geq0$ and $\eta,\tau>0$.  There is a finite angular-order
cutoff beyond which no nonzero counted input can have a pure
$\ell$ core field.  This remains true if arbitrary passive dielectric
layers outside the core make either or both polarization amplitudes
arbitrarily large.  In particular increasing $2\ell+1$ at a single
resonance does not by itself yield an unbounded counted subspace through
this fixed section.
\end{corollary}
\begin{proof}
The constant $B_\ell(s,t)$ tends to infinity, whereas
$\delta/(\eta\tau)$ is fixed.  All exterior amplification is already
included in the arbitrary coefficients $A,B$ and cancels in the ratio.
\end{proof}

\begin{remark}[Scope and the remaining constructive alternatives]
\label{cn:scope}
This is a genuine Maxwell exclusion for the stated spherical route,
not a uniform theorem for Conjecture~\ref{pr:channel-target}.
A coherent mixture of different angular orders has additional
interference terms with different radial factors; the proof does not
bound them.  Nor does it cover a dielectric structured throughout the
core, a spherical center moving away with angular order, or a plane
that does not pass through that center.  Such changes are allowed by
the unrestricted target and require a new argument.

The parameter regime $\delta\geq\eta\tau$ permits the return-flow
mechanism in Proposition~\ref{cn:odd-return}; it does not prove that an
arbitrarily large physical counted subspace exists.  A large bypass
tolerance, a choice of the full lossless scattering map as $T$, or a
single circulating high-$Q$ field is not by itself a quantifier
counterexample.  The throughput threshold still requires independent
physical inputs, and passage must hold for every coherent vector.
No admissible all-device counterexample for either threshold regime
is inferred here.
\end{remark}

\subsection{An actual local inverse-synthesis step for mixed angular orders}

The preceding obstruction does not make a mixed-order proposal
impossible.  Conversely, saying that radial layers can independently
tune every partial wave is an extra claim.  The following local version
of that claim can be proved for finite ordinary dielectrics, including
control of the interior amplitude rather than only the scattering phase.

\begin{theorem}[PROVED: local independent control of finitely many complex core amplitudes]
\label{cn:local-synthesis}
Fix $k_0>0$, $\varepsilon_+>1$, a radius $R_0>0$, an annulus
$0<r_-<r_+<R_0$, and distinct positive angular orders
$l_1,\ldots,l_m$.  Choose $1<n_b^2<\varepsilon_+$.
There exist $2m$ real smooth radial functions $h_j$ supported in
$(r_-,r_+)$ such that the finite scalar dielectric
\begin{equation}\label{cn:control-material}
 \varepsilon_\theta(r)=
 \begin{cases}
 n_b^2+\sum_{j=1}^{2m}\theta_jh_j(r),&r<R_0,\\
 1,&r>R_0
 \end{cases}
\end{equation}
has the following property for all sufficiently small real
$\theta\in\mathbb R^{2m}$.

For each $l=l_i$, fix the leading coefficient at the origin of the
regular real TE radial solution $u_l$ independently of $\theta$.
In the vacuum exterior write
\begin{equation}\label{cn:incoming-coefficient}
 u_l(r;\theta)=A_l(\theta)\,\xi_l^-(k_0r)
           +\overline{A_l(\theta)}\,\xi_l^+(k_0r),
 \qquad \xi_l^\pm(x)=x h_l^{(2/1)}(x).
\end{equation}
Here $\xi_l^-$ uses $h_l^{(2)}$ and $\xi_l^+$ uses
$h_l^{(1)}$.  Each $A_l$ is nonzero.  The map
\begin{equation}\label{cn:log-controls}
 \theta\longmapsto
 \left(\operatorname{Re}\log\frac{A_{l_i}(\theta)}{A_{l_i}(0)},
       \operatorname{Im}\log\frac{A_{l_i}(\theta)}{A_{l_i}(0)}
 \right)_{i=1}^m
\end{equation}
is a real-analytic local diffeomorphism at zero.  Thus every sufficiently
small prescribed change of all $m$ core gains and all $m$ scattering
phases is realized by a physical radial permittivity in the same fixed
interval $[1,\varepsilon_+]$.

Here is one explicit, retained neighborhood certificate.  Denote the
map in \eqref{cn:log-controls} by $\mathcal F$, let
$\sigma$ be the smallest singular value of $J=D\mathcal F(0)$,
and choose a parameter ball of radius $\rho$ on which
\begin{equation}\label{cn:control-radius}
 \rho\left(\sum_j\|h_j\|_\infty^2\right)^{1/2}
 <\min(n_b^2-1,\varepsilon_+-n_b^2),\qquad
 \sup_{|\theta|\leq\rho}\|D\mathcal F(\theta)-J\|
 \leq\sigma/2.
\end{equation}
The logarithms are taken on their continuous branches near zero.
Then every target vector of Euclidean norm less than
$\sigma\rho/2$ is attained in this parameter ball.
No lower bound on $\sigma\rho$ uniform in $m$ is asserted.
\end{theorem}

\begin{proof}
The TE equation, its interface conditions, and its regular normalization
are
\begin{equation}\label{cn:controlled-ode}
 u_l''+\left(k_0^2\varepsilon_\theta(r)
                 -\frac{l(l+1)}{r^2}\right)u_l=0,
 \qquad u_l,u_l'\text{ continuous}.
\end{equation}
Since every control vanishes near zero, the regular initial data can be
set at any fixed radius below $r_-$.  Standard successive integration
of this first-order system on a finite interval gives real-analytic
dependence on the finite parameters.  The two vacuum Hankel solutions
have Wronskian $2ik_0$, with the convention
$W(f,g)=fg'-f'g$.  The real nonzero regular solution has conjugate
exterior coefficients, so $A_l=0$ would force the solution to vanish
outside and then everywhere.  This proves its nonvanishing and the
local analytic logarithms.

Let $z_l(r)$ be the outgoing radial solution, equal to
$\xi_l^+(k_0r)$ outside $R_0$, continued inward through the base
material at $\theta=0$.  A real variation $h(r)$ changes the regular
solution by $\dot u_l$ satisfying
\begin{equation}\label{cn:variation-ode}
 \dot u_l''+\left(k_0^2\varepsilon_0-l(l+1)/r^2\right)\dot u_l
       =-k_0^2h u_l,
 \qquad \dot u_l=\dot u_l'=0\quad(r<r_-).
\end{equation}
Differentiate $W(\dot u_l,z_l)$, using the unperturbed equation for
$z_l$, and integrate.  Evaluation in the exterior gives the exact
first-variation formula
\begin{equation}\label{cn:amplitude-derivative}
 D\log A_l(0)[h]
   =\frac{k_0}{2iA_l(0)}\int_{r_-}^{r_+}
                    h(r)u_l(r)z_l(r)\,dr.
\end{equation}
This formula controls the complex incoming coefficient, not just the
phase of a unit-modulus scattering coefficient.

Put $\kappa=n_bk_0$, $\psi_l(x)=xj_l(x)$, and
$\zeta_l(x)=xy_l(x)$.  Inside the base sphere,
$u_l=c_l\psi_l(\kappa r)$ with real $c_l\ne0$, while
\begin{equation}\label{cn:outgoing-interior}
 z_l=\alpha_l\psi_l(\kappa r)+\beta_l\zeta_l(\kappa r).
\end{equation}
The real Wronskian of $\psi_l(\kappa r)$ and
$\zeta_l(\kappa r)$ is $\kappa$.  In the vacuum exterior,
$W(z_l,\overline{z_l})=-2ik_0$; consequently
\begin{equation}\label{cn:independent-coefficients}
 \operatorname{Im}(\alpha_l\overline{\beta_l})
       =-k_0/\kappa\ne0.
\end{equation}
Thus the real and imaginary parts of the kernel in
\eqref{cn:amplitude-derivative} are obtained by an invertible real
two-by-two change of basis from
\begin{equation}\label{cn:product-kernels}
 \psi_l(\kappa r)^2,\qquad
 \psi_l(\kappa r)\zeta_l(\kappa r).
\end{equation}

We now prove independence of all $2m$ real functions in
\eqref{cn:product-kernels}.  Their expressions extend analytically
through zero: the first functions are even and the second odd, since
the zero of $\psi_l$ cancels the pole of $\zeta_l$.
An identity on the prescribed annulus therefore extends to all real
arguments and separates into its even and odd parts.
The squares are independent because their first nonzero Taylor
orders $2l_i+2$ are distinct.

For the odd products use the exact finite Hankel expressions
\begin{equation}\label{cn:finite-hankel}
 \xi_l^+(x)=(-i)^{l+1}e^{ix}p_l(x),\qquad
 \xi_l^-(x)=i^{l+1}e^{-ix}\widetilde p_l(x),
\end{equation}
where $p_l,\widetilde p_l$ are polynomials of degree $l$ in $1/x$
with nonzero leading coefficients.  These follow directly by the
spherical-Bessel recurrence starting from the sine and cosine
solutions at $l=0$; they are also displayed in
\eqref{ms:hankel-polynomials}.
Since
\begin{equation}\label{cn:odd-product-hankel}
 \psi_l\zeta_l=
 \frac{(\xi_l^+)^2-(\xi_l^-)^2}{4i},
\end{equation}
a nontrivial linear relation among the odd products would imply
$e^{4ix}P(1/x)=Q(1/x)$ for two polynomials $P,Q$.
If $P$ is nonzero, the quotient $Q(1/x)/P(1/x)$ has either a finite
limit, zero limit, or diverging modulus as $x\to+\infty$; none equals
$e^{4ix}$.  Thus $P=Q=0$.  The coefficient of the highest remaining
power $x^{-2l_i}$ in $P$ then forces the coefficient of the largest
$l_i$ to vanish.  Descending induction forces all coefficients to
vanish.  This proves the claimed independence.

For any $2m$ independent continuous real functions on an interval,
there exist $2m$ evaluation points whose value matrix is invertible.
Indeed, if the span of all evaluation vectors were proper, a nonzero
linear combination of the functions would vanish everywhere.
Choose such points for the real derivative kernels in
\eqref{cn:amplitude-derivative}.  Real smooth unit-integral bumps
$h_j$ in sufficiently small, disjoint neighborhoods of those points
make the derivative matrix as close as desired to that invertible
evaluation matrix.  Hence $J$ is invertible and $\sigma>0$.
The inverse-function theorem proves the qualitative local assertion.

For the retained certificate, write
$\mathcal F(\theta)=J\theta+\mathcal R(\theta)$.
The second condition in \eqref{cn:control-radius} makes
$\mathcal R$ Lipschitz with constant $\sigma/2$ on the parameter
ball and $\mathcal R(0)=0$.  For a target $y$ with
$|y|<\sigma\rho/2$, the map
$\theta\mapsto J^{-1}(y-\mathcal R(\theta))$
is a contraction of that ball into itself.  Its fixed point realizes
$y$.  The first condition in \eqref{cn:control-radius} keeps the
permittivity strictly within its prescribed bounds.

Finally, for the angular normalization
$\int_{S^2}|X_{lm}|^2\,dS=1$, an incoming TE radial amplitude
$A$ carries physical power $|A|^2/(2Z_0)$ in the vacuum exterior.
Thus the field produced by unit incoming power has regular-core
amplitude $\sqrt{2Z_0}/A_l$ in the normalization above, while its
outgoing-to-incoming coefficient is
$\overline{A_l}/A_l$.  The real and imaginary logarithms controlled
in \eqref{cn:log-controls} are precisely the logarithmic core gain
with a minus sign, and minus one half the scattering phase.
The vacuum power convention does not change with the controls.
\end{proof}

\begin{remark}[The local synthesis range is not an arbitrary matching network]
\label{cn:local-synthesis-scope}
Theorem~\ref{cn:local-synthesis} supplies actual radial Maxwell
material controls for a finite set of angular orders.  It neither
makes $A_l$ arbitrarily small simultaneously nor estimates the
absolute flux on the whole section plane.  The certified target
neighborhood may shrink with the number and size of the angular
orders.  Iterating local invertibility without a continuation
estimate would not prove global reachability within the fixed
permittivity bounds.  Independent tuning of finitely many scattering
phases alone would be still weaker, because it would not control
the gain represented by $|A_l|^{-1}$.
\end{remark}

\begin{corollary}[PROVED: local control of both polarizations at separated orders]
\label{cn:both-synthesis}
If the selected positive angular orders additionally satisfy
$|l_i-l_j|\geq3$ for $i\ne j$, there exist $4m$ smooth real radial
controls in the same prescribed annulus that independently vary,
locally, the complex regular-core incoming coefficients of both
TE and TM polarizations at all selected orders.  All material and
physical normalization conclusions of
Theorem~\ref{cn:local-synthesis} remain valid, with a $4m$-dimensional
Jacobian and its own retained neighborhood constants.
\end{corollary}
\begin{proof}
For the TM polarization choose the magnetic radial amplitude $u$.
Its real equation and transmission conditions are
\begin{equation}\label{cn:tm-ode}
 (\varepsilon^{-1}u')'
       +\left(k_0^2-\frac{l(l+1)}{\varepsilon r^2}\right)u=0,
 \qquad u,\ \varepsilon^{-1}u'\text{ continuous}.
\end{equation}
It has the same regular/irregular basis inside a constant-index core
and the same vacuum radial Hankel basis outside.  If $z_l$ is its
continued outgoing solution, differentiating the weighted Wronskian
$\varepsilon^{-1}W(\dot u_l,z_l)$ and integrating by parts gives
\begin{equation}\label{cn:tm-amplitude-derivative}
 D\log A_l^{\rm TM}(0)[h]
  =\frac{1}{2ik_0 A_l^{\rm TM}(0)n_b^4}
       \int h(r)\left(u_l'z_l'
                  +\frac{l(l+1)}{r^2}u_lz_l\right)\,dr.
\end{equation}
Indeed $\dot p=-h/n_b^4$ for $p=\varepsilon^{-1}$,
and the integrated weighted-Wronskian derivative is
$-\int\dot p(u_l'z_l'+l(l+1)u_lz_l/r^2)$.
The endpoint terms vanish since the control has compact support
inside the annulus.  The exterior weighted Wronskian is again
$2ik_0\dot A_l^{\rm TM}$.

Inside the sphere write $u_l=c_l\psi_l(\kappa r)$ and
$z_l=\alpha_l\psi_l(\kappa r)+\beta_l\zeta_l(\kappa r)$.
The weighted outgoing Wronskian gives
$\operatorname{Im}(\alpha_l\overline{\beta_l})=-n_b\ne0$.
Consequently the real and imaginary derivative kernels are an
invertible real change of basis from
\begin{equation}\label{cn:tm-kernels}
 (\psi_l')^2+\frac{l(l+1)}{x^2}\psi_l^2,
 \qquad
 \psi_l'\zeta_l'+\frac{l(l+1)}{x^2}\psi_l\zeta_l,
 \quad x=\kappa r.
\end{equation}
The common factor $\kappa^2$ is immaterial for independence.

The elementary recurrences
$\psi_{l-1}=\psi_l'+l\psi_l/x$ and
$\psi_{l+1}=(l+1)\psi_l/x-\psi_l'$, and the identical
recurrences for $\zeta_l$, imply
\begin{align}
 (\psi_l')^2+\frac{l(l+1)}{x^2}\psi_l^2
   &=\frac{l+1}{2l+1}\psi_{l-1}^2
       +\frac{l}{2l+1}\psi_{l+1}^2,\notag\\
 \psi_l'\zeta_l'+\frac{l(l+1)}{x^2}\psi_l\zeta_l
   &=\frac{l+1}{2l+1}\psi_{l-1}\zeta_{l-1}
       +\frac{l}{2l+1}\psi_{l+1}\zeta_{l+1}.
 \label{cn:tm-adjacent-identities}
\end{align}
The independence proof for the product kernels in the preceding
theorem also includes order zero.  The triples
$\{l_i-1,l_i,l_i+1\}$ are disjoint by hypothesis.  The TE pair
uses the middle index of each triple; the TM pair is a nonzero
combination of its two unused neighboring indices, separately in
the even and odd kernel families.  All $4m$ derivative functions
are therefore independent.  The same evaluation-point, smooth-bump,
and inverse-function argument now proves the corollary.  For the
magnetic normalization of TM waves, incoming power is
$Z_0|A|^2/2$, another fixed vacuum factor independent of the controls.
\end{proof}

\begin{proposition}[PROVED: fixed-annulus controls have a finite gain ceiling]
\label{cn:fixed-control-ceiling}
Keep $r_->0$, $R_0$, the interior material on $r<r_-$, and an angular
order $l$ fixed.  Let the material vary arbitrarily on
$r_-<r<R_0$ within $1\leq\varepsilon\leq\varepsilon_+$,
with finitely many smooth pieces, and let it be vacuum beyond $R_0$.
For either polarization, the incoming coefficient $A_l$ of the same
origin-normalized regular solution obeys a strictly positive lower
bound independent of the number and thickness of radial layers.

Here is an explicit bound.  For TE let
$v=(u,u'/k_0)^T$; for TM let
$v=(u,\varepsilon^{-1}u'/k_0)^T$.  Set
\begin{align}
 C_l&=k_0(R_0-r_-)
       \max\left(\varepsilon_+,
               1+\frac{l(l+1)}{k_0^2r_-^2}\right),\notag\\
 V_l&=\begin{pmatrix}
 \operatorname{Re}\xi_l^-(k_0R_0)&
             -\operatorname{Im}\xi_l^-(k_0R_0)\\
 \operatorname{Re}(\xi_l^-)'(k_0R_0)&
             -\operatorname{Im}(\xi_l^-)'(k_0R_0)
 \end{pmatrix}.
 \label{cn:ceiling-data}
\end{align}
The primes in $V_l$ are argument derivatives.  Then
\begin{equation}\label{cn:amplitude-ceiling}
 |A_l|\geq
       \frac{e^{-C_l}|v(r_-)|}{2\|V_l\|}>0.
\end{equation}
\end{proposition}
\begin{proof}
The states are continuous across every radial interface.  Away from
interfaces their first-order systems are respectively
\begin{equation}\label{cn:state-generators}
 v'=k_0\begin{pmatrix}0&1\\
 l(l+1)/(k_0^2r^2)-\varepsilon&0\end{pmatrix}v,
 \qquad
 v'=k_0\begin{pmatrix}0&\varepsilon\\
 l(l+1)/(k_0^2\varepsilon r^2)-1&0\end{pmatrix}v.
\end{equation}
For a real matrix with off-diagonal entries $a,b$ and zero diagonal,
the Euclidean operator norm is $\max(|a|,|b|)$.  Both generators
are therefore bounded by $C_l/(R_0-r_-)$ on the interval.
Gronwall applied also to reverse propagation gives
$|v(R_0)|\geq e^{-C_l}|v(r_-)|$.
The initial state is fixed and nonzero by uniqueness of the regular
radial solution.  On the vacuum side of the final interface,
\eqref{cn:incoming-coefficient} gives
$v(R_0)=2V_l(\operatorname{Re}A_l,\operatorname{Im}A_l)^T$.
The Hankel Wronskian makes $V_l$ invertible, and its norm is finite.
This proves \eqref{cn:amplitude-ceiling}.
\end{proof}

\begin{remark}[Focused global continuation attempt and its exact limits]
\label{cn:global-attempt}
The local controls can leave every selected scattering phase fixed
while changing the selected core gains: prescribe zero imaginary
coordinates in \eqref{cn:log-controls}.  Thus phase matching by itself
does not determine the internal response.  But the controls cannot
be iterated inside their fixed bounded annulus to obtain arbitrarily
large gains; Proposition~\ref{cn:fixed-control-ceiling} gives an
explicit obstruction to that proposed iteration, even with
arbitrarily many increasingly thin layers.

Allowing the controlled radius to increase removes this particular
ceiling, as the frozen unrestricted target permits.  It also adds an
uncontrolled portion of the physical section plane.  Endpoint radial
transfer data $A_l$ do not estimate its absolute normal flux.  The
single-order estimate in Theorem~\ref{cn:radial-bypass} gives a
concrete case in which the exterior plane cost grows rather than
disappears under amplification.  No reachability proof with increasing
support and simultaneous whole-plane bypass control has been
obtained here.  These are the exact two missing estimates in this
mixed-order construction, not unspecified corrections to a completed
counterexample.  The unrestricted finite-channel statement is still
unresolved by this route.
\end{remark}
